\documentclass[10pt,a4paper]{amsart}
\usepackage{amsmath,amssymb,amsthm,graphicx}
\usepackage{verbatim}
\newenvironment{reference}{\comment}{\endcomment}
%\newenvironment{reference}{}{}
\newenvironment{slogan}{\comment}{\endcomment}
\newenvironment{history}{\comment}{\endcomment}

% For commutative diagrams we use Xy-pic
\usepackage[all]{xy}

% We use 2cell for 2-commutative diagrams.
\xyoption{2cell}
% Theorem environments.
%
\theoremstyle{plain}
\newtheorem{theorem}[subsection]{Theorem}
\newtheorem{proposition}[subsection]{Proposition}
\newtheorem{lemma}[subsection]{Lemma}

\theoremstyle{definition}
\newtheorem{definition}[subsection]{Definition}
\newtheorem{example}[subsection]{Example}
\newtheorem{exercise}[subsection]{Exercise}
\newtheorem{situation}[subsection]{Situation}

\theoremstyle{remark}
\newtheorem{remark}[subsection]{Remark}
\newtheorem{remarks}[subsection]{Remarks}

\numberwithin{equation}{subsection}

% Macros
%
\def\lim{\mathop{\mathrm{lim}}\nolimits}
\def\colim{\mathop{\mathrm{colim}}\nolimits}
\def\Spec{\mathop{\mathrm{Spec}}}
\def\Hom{\mathop{\mathrm{Hom}}\nolimits}
\def\Ext{\mathop{\mathrm{Ext}}\nolimits}
\def\SheafHom{\mathop{\mathcal{H}\!\mathit{om}}\nolimits}
\def\SheafExt{\mathop{\mathcal{E}\!\mathit{xt}}\nolimits}
\def\Sch{\mathit{Sch}}
\def\Mor{\mathop{\mathrm{Mor}}\nolimits}
\def\Ob{\mathop{\mathrm{Ob}}\nolimits}
\def\Sh{\mathop{\mathit{Sh}}\nolimits}
\def\NL{\mathop{N\!L}\nolimits}
\def\CH{\mathop{\mathrm{CH}}\nolimits}
\def\proetale{{pro\text{-}\acute{e}tale}}
\def\etale{{\acute{e}tale}}
\def\QCoh{\mathit{QCoh}}
\def\Ker{\mathop{\mathrm{Ker}}}
\def\Im{\mathop{\mathrm{Im}}}
\def\Coker{\mathop{\mathrm{Coker}}}
\def\Coim{\mathop{\mathrm{Coim}}}

% Boxtimes
%
\DeclareMathSymbol{\boxtimes}{\mathbin}{AMSa}{"02}

%
% Macros for moduli stacks/spaces
%
\def\QCohstack{\mathcal{QC}\!\mathit{oh}}
\def\Cohstack{\mathcal{C}\!\mathit{oh}}
\def\Spacesstack{\mathcal{S}\!\mathit{paces}}
\def\Quotfunctor{\mathrm{Quot}}
\def\Hilbfunctor{\mathrm{Hilb}}
\def\Curvesstack{\mathcal{C}\!\mathit{urves}}
\def\Polarizedstack{\mathcal{P}\!\mathit{olarized}}
\def\Complexesstack{\mathcal{C}\!\mathit{omplexes}}
% \Pic is the operator that assigns to X its picard group, usage \Pic(X)
% \Picardstack_{X/B} denotes the Picard stack of X over B
% \Picardfunctor_{X/B} denotes the Picard functor of X over B
\def\Pic{\mathop{\mathrm{Pic}}\nolimits}
\def\Picardstack{\mathcal{P}\!\mathit{ic}}
\def\Picardfunctor{\mathrm{Pic}}
\def\Deformationcategory{\mathcal{D}\!\mathit{ef}}
\usepackage{cleveref}
\crefname{originalreference}{original reference}{original references}
\Crefname{originalreference}{Original reference}{Original references}
\makeatletter
\newlabel{situation-setup}{{02QL}{source}{}{sourcepointer-1}{}}
\newlabel{situation-setup@cref}{{[originalreference][0][]02QL}{[1][0][]source}{}{}{}}
\makeatother
\crefname{originaldefinition}{original definition}{original definitions}
\Crefname{originaldefinition}{Original definition}{Original definitions}
\makeatletter
\newlabel{definition-gysin-homomorphism}{{02T8}{source}{}{sourcepointer-2}{}}
\newlabel{definition-gysin-homomorphism@cref}{{[originaldefinition][0][]02T8}{[1][0][]source}{}{}{}}
\makeatother
\crefname{originallemma}{original lemma}{original lemmas}
\Crefname{originallemma}{Original lemma}{Original lemmas}
\makeatletter
\newlabel{lemma-gysin-bivariant}{{0B79}{source}{}{sourcepointer-3}{}}
\newlabel{lemma-gysin-bivariant@cref}{{[originallemma][0][]0B79}{[1][0][]source}{}{}{}}
\makeatother
\crefname{originallemma}{original lemma}{original lemmas}
\Crefname{originallemma}{Original lemma}{Original lemmas}
\makeatletter
\newlabel{lemma-bivariant-zero}{{02UC}{source}{}{sourcepointer-4}{}}
\newlabel{lemma-bivariant-zero@cref}{{[originallemma][0][]02UC}{[1][0][]source}{}{}{}}
\makeatother
\crefname{originallemma}{original lemma}{original lemmas}
\Crefname{originallemma}{Original lemma}{Original lemmas}
\makeatletter
\newlabel{lemma-gysin-fundamental}{{0FBM}{source}{}{sourcepointer-5}{}}
\newlabel{lemma-gysin-fundamental@cref}{{[originallemma][0][]0FBM}{[1][0][]source}{}{}{}}
\makeatother
\title{The Gysin homomorphism of an effective Cartier divisor agrees with the class of its virtual normal sheaf}
\author{The Stacks Project Authors}
\date{stacks-project:chow.tex:local-snapshot-20261012; GNU FDL 1.2 or later}
\begin{document}
\maketitle
\noindent\textit{Editorial note.} All delivered mathematical text is the author's own, copied from the pinned Stacks Project chapter file chow.tex; no statement or proof step was rewritten, reordered or removed.\par\medskip
\noindent\textit{Editorial note.} Preamble provenance: the pinned chapter file opens with an input of the shared preamble and carries no preamble text of its own, so the project's single Stacks-wide preamble staged at sources/stacks-chow/preamble.tex is excerpted verbatim as two spans, lines 14-18 (Xy-pic setup) and lines 154-217 (the native theorem declarations and mathematical macros). Nothing was added, removed or redefined; the wrapper class and the amsmath/amssymb/amsthm/graphicx base packages are supplied by the assembler.\par\medskip
\noindent\textit{Editorial note.} The counted statement and proof contain no cross reference, no citation and no included file, so the delivered body is self contained.\par\medskip
\noindent\textit{Editorial note.} Canonical location of the counted proposition: Stacks tag 0FBN (sources/stacks-chow/pinned-original-source.tex line 12174).\par\medskip
\noindent\textit{Allowed standard prerequisites (editorial).}\begin{itemize}
\item The Gysin class attached to a divisor and the virtual normal sheaf. Same chapter, Definition labelled definition-gysin-homomorphism and Lemma labelled lemma-gysin-bivariant, invoked in the statement. Fixes the two classes being compared.
\item Testing identities of bivariant classes on an integral base. Same chapter, Lemma labelled lemma-bivariant-zero, cited at proof lines 12185-12187. Reduces to X integral and the identity i\textasciicircum{}*[X] = c cap [X].
\item The fundamental computation for the class of a virtual normal sheaf. Same chapter, Lemma labelled lemma-gysin-fundamental, cited at proof lines 12195-12198 and 12200-12203. Evaluates the class in both cases.
\end{itemize}\medskip
\ifdefined\hypertarget\hypertarget{sourcepointer-1}{}\fi
\noindent\textit{Source reference.} Original reference 02QL: sources/stacks-chow/pinned-original-source.tex, source lines 1372--1372. Same-chapter pointer into the author environment situation at source line 1372; invoked at proof line None of the counted proof. It is retained only as a navigation pointer; every load-bearing step is inside the delivered span.\par\medskip
\ifdefined\hypertarget\hypertarget{sourcepointer-2}{}\fi
\noindent\textit{Source reference.} Original definition 02T8: sources/stacks-chow/pinned-original-source.tex, source lines 5069--5069. Same-chapter pointer into the author environment definition at source line 5069; invoked at proof line 12194 of the counted proof. It is retained only as a navigation pointer; every load-bearing step is inside the delivered span.\par\medskip
\ifdefined\hypertarget\hypertarget{sourcepointer-3}{}\fi
\noindent\textit{Source reference.} Original lemma 0B79: sources/stacks-chow/pinned-original-source.tex, source lines 6149--6149. Same-chapter pointer into the author environment lemma at source line 6149; invoked at proof line None of the counted proof. It is retained only as a navigation pointer; every load-bearing step is inside the delivered span.\par\medskip
\ifdefined\hypertarget\hypertarget{sourcepointer-4}{}\fi
\noindent\textit{Source reference.} Original lemma 02UC: sources/stacks-chow/pinned-original-source.tex, source lines 6585--6585. Same-chapter pointer into the author environment lemma at source line 6585; invoked at proof line 12186 of the counted proof. It is retained only as a navigation pointer; every load-bearing step is inside the delivered span.\par\medskip
\ifdefined\hypertarget\hypertarget{sourcepointer-5}{}\fi
\noindent\textit{Source reference.} Original lemma 0FBM: sources/stacks-chow/pinned-original-source.tex, source lines 12076--12076. Same-chapter pointer into the author environment lemma at source line 12076; invoked at proof line 12193 of the counted proof. It is retained only as a navigation pointer; every load-bearing step is inside the delivered span.\par\medskip
\begin{lemma}
\label{lemma-gysin-agrees}
Let $(S, \delta)$ be as in Situation \ref{situation-setup}. Let $X$ be a scheme
locally of finite type over $S$. Let $(\mathcal{L}, s, i : D \to X)$
be a triple as in Definition \ref{definition-gysin-homomorphism}.
The Gysin homomorphism $i^*$ viewed as an element of $A^1(D \to X)$
(see Lemma \ref{lemma-gysin-bivariant}) is the same as the bivariant class
$c(D \to X, \mathcal{N}) \in A^1(D \to X)$
constructed using $\mathcal{N} = i^*\mathcal{L}$
viewed as a virtual normal sheaf for $D$ in $X$.
\end{lemma}

\begin{proof}
We will use the criterion of Lemma \ref{lemma-bivariant-zero}.
Thus we may assume that $X$ is an integral scheme and
we have to show that $i^*[X]$ is equal to $c \cap [X]$.
Let $n = \dim_\delta(X)$. As usual, there are two cases.

\medskip\noindent
If $X = D$, then we see that both classes are represented by
$c_1(\mathcal{N}) \cap [X]_n$. See Lemma \ref{lemma-gysin-fundamental}
and Definition \ref{definition-gysin-homomorphism}.

\medskip\noindent
If $D \not = X$, then $D \to X$ is an effective Cartier divisor
and in particular a regular closed immersion of codimension $1$.
Again by Lemma \ref{lemma-gysin-fundamental} we conclude
$c(D \to X, \mathcal{N}) \cap [X]_n = [D]_{n - 1}$. The same
is true by definition for the Gysin homomorphism and we conclude
once again.
\end{proof}

\end{document}
